% ==============================================================================
% SLIDES - SEMANA 15: IMPLANTAÇÃO DA GOVERNANÇA, FCS & GESTÃO DA MUDANÇA
% ==============================================================================

% ------------------------------------------------------------------------------
% 1. SLIDE DE CAPA
% ------------------------------------------------------------------------------
\senaicover
  {Implantação da Governança de TI}
  {Roadmap em Ondas, Fatores Críticos de Sucesso, Mudança (ADKAR) e Orçamento}
  {Unidade Curricular: Governança de TI}
  {Prof. André Souza}

% ------------------------------------------------------------------------------
% 2. VISÃO GERAL / AGENDA
% ------------------------------------------------------------------------------
\begin{senaiframe}{Visão Geral \& Agenda}{Os 4 Eixos da Aula}

  \begin{columns}[t]
    \begin{column}{0.48\textwidth}
      \begin{tcolorbox}[senaibluecard, title={1. Roadmap em Ondas}]
        \begin{itemize}
          \item Onda 1 (Curto Prazo): Quick Wins.
          \item Onda 2 (Médio) e Onda 3 (Longo Prazo).
        \end{itemize}
      \end{tcolorbox}
      
      \vspace{4pt}
      
      \begin{tcolorbox}[senaiambercard, title={3. Gestão da Mudança}]
        \begin{itemize}
          \item Modelo ADKAR (Prosci).
          \item Vencendo a resistência cultural.
        \end{itemize}
      \end{tcolorbox}
    \end{column}

    \begin{column}{0.48\textwidth}
      \begin{tcolorbox}[senairedcard, title={2. Fatores Críticos}]
        \begin{itemize}
          \item Patrocínio Executivo (Tone at the Top).
          \item Foco nos negócios e vitórias rápidas.
        \end{itemize}
      \end{tcolorbox}
      
      \vspace{4pt}
      
      \begin{tcolorbox}[senaigreencard, title={4. Capítulo 9 do PDGTI}]
        \begin{itemize}
          \item Roadmap e Investimentos (CAPEX/OPEX).
          \item Matriz de Geração de Valor.
        \end{itemize}
      \end{tcolorbox}
    \end{column}
  \end{columns}

\end{senaiframe}

% ------------------------------------------------------------------------------
% 3. ROADMAP EM ONDAS
% ------------------------------------------------------------------------------
\begin{senaiframe}{Módulo 1 • Roteiro}{Implantação Estruturada em 3 Ondas}

  \begin{columns}[t]
    \begin{column}{0.48\textwidth}
      \begin{tcolorbox}[senaibluecard, title={Onda 1 — Curto Prazo (0-6m)}]
        \begin{itemize}
          \item Contenção de danos, inventário inicial, controles de acesso e criação do Comitê Estratégico de TI.
        \end{itemize}
      \end{tcolorbox}
    \end{column}

    \begin{column}{0.48\textwidth}
      \begin{tcolorbox}[senaicard, title={Ondas 2 \& 3 — Médio/Longo}]
        \begin{itemize}
          \item \textbf{Onda 2 (6-18m):} Adoção ITIL/COBIT, formalização do PETI e monitoramento.
          \item \textbf{Onda 3 (18-36m):} Maturidade CMMI/MPS-SW e automação avançada.
        \end{itemize}
      \end{tcolorbox}
    \end{column}
  \end{columns}

\end{senaiframe}

% ------------------------------------------------------------------------------
% 4. GESTÃO DA MUDANÇA (ADKAR) & ORÇAMENTO
% ------------------------------------------------------------------------------
\begin{senaiframe}{Módulo 2 • Pessoas \& Recursos}{O Modelo ADKAR \& CAPEX vs. OPEX}

  \begin{columns}[t]
    \begin{column}{0.48\textwidth}
      \begin{tcolorbox}[senaiambercard, title={Modelo ADKAR (Mudança)}]
        \begin{itemize}
          \item \textbf{Awareness:} Conscientização do motivo.
          \item \textbf{Desire:} Desejo pessoal de apoiar.
          \item \textbf{Knowledge \& Ability:} Saber e praticar.
          \item \textbf{Reinforcement:} Sustentação contínua.
        \end{itemize}
      \end{tcolorbox}
    \end{column}

    \begin{column}{0.48\textwidth}
      \begin{tcolorbox}[senairedcard, title={CAPEX vs. OPEX}]
        \begin{itemize}
          \item \textbf{CAPEX:} Investimento em ativos físicos (hardware, data center).
          \item \textbf{OPEX:} Despesas operacionais (SaaS, nuvem por assinatura).
        \end{itemize}
      \end{tcolorbox}
    \end{column}
  \end{columns}

\end{senaiframe}

% ------------------------------------------------------------------------------
% 5. REFERÊNCIAS BIBLIOGRÁFICAS
% ------------------------------------------------------------------------------
\begin{senaiframe}{Referências Bibliográficas}{Fontes \& Leitura Recomendada}

  \begin{tcolorbox}[senaicard, title={Obras de Referência}]
    \begin{itemize}
      \item \textbf{FERNANDES, Aguinaldo Aragon; ABREU, Vladimir Ferraz.} *Implantando a governança de TI: da estratégia à gestão de processos e serviços.* 4. ed. Rio de Janeiro: Brasport, 2014.
      \item \textbf{HIATT, Jeffrey M.} *ADKAR: a model for change in business, government and our community.* Prosci Learning Center Publications, 2006.
    \end{itemize}
  \end{tcolorbox}

\end{senaiframe}

% ------------------------------------------------------------------------------
% 6. APLICAÇÃO NO PDGTI
% ------------------------------------------------------------------------------
\begin{senaiframe}{Aplicação no PDGTI}{Desenvolvimento do Capítulo 9}

  \vspace{0.2cm}
  \begin{tcolorbox}[senaigreencard, title={Capítulo 9 – Roadmap de Implantação \& Benefícios}]
    \begin{itemize}
      \item \textbf{9.1 Cronograma em Ondas:} Tabela ou Gantt com as 3 ondas de implantação.
      \item \textbf{9.2 Estimativa Orçamentária:} Projeção de investimentos em CAPEX e OPEX.
      \item \textbf{9.3 Fatores Críticos de Sucesso:} Mapear no mínimo 4 FCS específicos.
      \item \textbf{9.4 Matriz de Benefícios:} Demonstrar a geração de valor para o negócio.
    \end{itemize}
  \end{tcolorbox}

\end{senaiframe}
